\chapter*{Sommario}
Questa tesi di dottorato consiste principalmente in una presentazione organizzata e coerente di due articoli di ricerca \cite{Amb:fil-fuk} e \cite{ABC26} co-firmati dall'autore. Il lavoro è dedicato alla definizione della categoria di Fukaya derivata persistente, un raffinamento della nozione di categoria di Fukaya derivata —ormai classica della topologia simplettica — e allo studio di alcune delle sue proprietà.

La tesi è suddivisa in tre parti. Nella prima parte vengono introdotti gli strumenti algebrici necessari per la definitione e lo studio delle categorie di Fukaya derivate persistenti, in particolare le categorie $A_\infty$ filtrate, le categorie triangolate persistenti (TPCs) e la nozione astratta di approssimabilità metrica categorica, interpretabile sia come una generalizzazione della nozione metrica di spazio totalmente limitato, sia come una versione persistente del concetto di generazione triangolare nelle categorie triangolate.\\
Nella seconda parte viene sviluppato un modello di categoria di Fukaya filtrata per alcune classi di varietà simplettiche. Questo passaggio preliminare è necessario per la definizione di categorie di Fukaya derivate persistenti. Gli strumenti principali sono: (i) un modello \say{cluster} per la categoria di Fukaya e (ii) la costruzione di una classe adeguata di perturbazioni hamiltoniane.\\
Nella terza parte si dimostra che alcune classi di sottovarietà lagrangiane, dotate della metrica spettrale lagrangiana, sono categoricamente approssimabili in opportune categorie di Fukaya derivate persistenti. Si mostra inoltre che, in alcuni casi, tali classi sono approssimabili ma non totalmente limitate.

\chapter*{Abstract}

This doctoral thesis consists in an organised and coherent presentation of two articles \cite{Amb:fil-fuk} and \cite{ABC26} co-authored by the present author. The thesis is devoted to the definition of the persistent derived Fukaya category, which refines the classical notion of the derived Fukaya category in symplectic topology, and to the study of some of its main features.

The thesis consists of three parts. In the first part we introduce all the necessary algebraic material needed to discuss persistence derived Fukaya categories and categorical metric approximability, that is, filtered $A_\infty$-categories, TPCs and the abstract notion of categorical metric approximability, which can be seen as a weaker version of total boundedness as well as a persistence version of triangular generation in triangulated categories. \\ In the second part we develop a model for a filtered Fukaya category of certain symplectic manifolds. This is a necessary step in order to define TPC refinements of derived Fukaya categories. The main tools are: (i) a cluster model for the Fukaya category and (ii) the construction of an adequate class of Hamiltonian perturbation data.\\
In the third part we show that certain classes of Lagrangian submanifolds endowed with the Lagrangian spectral metric are categorically metric approximable in appropriate persistence derived Fukaya categories. We show that, in some cases, these classes are approximable but not totally bounded.
